\documentclass[11pt,a4paper]{article}
\usepackage[utf8]{inputenc}
\usepackage{ctex} % 完美支持中文
\usepackage{amsmath, amssymb, amsthm}
\usepackage{geometry}
\usepackage{xcolor}
\usepackage{tcolorbox} % 用于打造高颜值边框

% 设定优雅的页面边距
\geometry{left=2.3cm,right=2.3cm,top=2.5cm,bottom=2.5cm}

% 自定义高颜值卡片样式 (用于解答题)
\tcbset{
	examstyle/.style={
		colback=gray!5, % 卡片背景色：极淡灰
		colframe=navyblue!70!black, % 边框颜色
		arc=3mm, % 圆角
		boxrule=0.5pt, % 边框线条粗细
		left=4mm, right=4mm, top=3mm, bottom=3mm,
		before=\vspace{10pt}, after=\vspace{10pt}
	}
}
\definecolor{navyblue}{RGB}{30, 60, 90}

% 重新定义微分算子，使其更符合出版级学术规范
\newcommand{\dd}{\mathrm{d}}

\title{\vspace{-1cm}\Huge\textbf{华东师范大学 2026 年数学分析试卷}\\[0.3em]\large --- 考研真题与深度解析标准版 ---}
\author{\textbf{数学分析专家组}}
\date{\today}

\begin{document}
	
	\maketitle
	
	\begin{center}
		\rule{\textwidth}{0.8pt}
	\end{center}
	
	\section{一、 填空题（每题 8 分，共 56 分）}
	
	\begin{enumerate}
		\item 计算 $\sum_{n=1}^\infty \frac{2n - 1}{2^n} = \underline{\makebox[2cm]{3}}$。
		
		\item 计算 $\lim_{x \to 0^+} \sqrt[x]{\sec \sqrt{2x}} = \underline{\makebox[2cm]{e}}$。
		
		\item 计算 $\int_0^1 \int_0^{\sqrt{1-x^2}} \int_{\sqrt{x^2+y^2}}^{\sqrt{2-x^2-y^2}} xy \,\dd z\dd y\dd x = \underline{\makebox[2cm]{$\frac{4\sqrt{2}-5}{30}$}}$。
		
		\item 设 $|a| \neq 1$，计算 $\int_0{\frac{\pi}{2}} \frac{\arctan(a \tan x)}{\tan x} \,\dd x = \underline{\makebox[3.5cm]{$\frac{\pi}{2} \text{sgn}(a) \ln(1+|a|)$}}$。
		
		\item 求和函数 $\sum_{n=1}^\infty \frac{(-1)^{n-1}}{n(n+1)}(x^2 - 1)^{2n} = \underline{\makebox[6cm]{$\left(1 - \frac{1}{(x^2-1)^2}\right)\ln\left(1+(x^2-1)^2\right) + 1$}}$。
		
		\item 计算第二类曲面积分 $\iint_\Sigma xz^2 \,\dd y\dd z + \frac{1}{y} \,\dd z\dd x + x^2z \,\dd x\dd y = \underline{\makebox[3cm]{$\frac{3}{2}\pi - \pi\ln 2$}}$，其中 $\Sigma$ 为抛物面 $y = x^2+z^2$ 被平面 $y=1, y=2$ 所截部分，取外侧。
		
		\item 设 $0 < x < 2\pi, a \neq 0$，则 $\frac{1}{2a} + \sum_{n=1}^\infty \frac{a \cos nx - n \sin nx}{n^2 + a^2} = \underline{\makebox[4cm]{$\frac{\pi e^{a(\pi-x)}}{2\sinh(a\pi)}$}}$。
	\end{enumerate}
	
	\newpage
	\section{二、 解答题（需提供以下各题的详细解答过程）}
	
	% 第8题卡片
	\begin{tcolorbox}[examstyle]
		\subsection*{第 8 题：方程根的性态与数列收敛（15 分）}
		\textbf{题目：} 
		(1) 证明：对任意正整数 $n$，方程 $x^3 + 3x + \frac{1}{n} = 0$ 有且仅有一个实数解。 \\
		(2) 设 $x_n \, (n = 1,2,\dots)$ 为 (1) 中方程的实根，证明：数列 $\{x_n\}$ 收敛。
		
		\begin{proof}[\textbf{证明}]
			(1) 令 $f_n(x) = x^3 + 3x + \frac{1}{n}$。对其求导得 $f_n'(x) = 3x^2 + 3 = 3(x^2+1) > 0$ 恒成立。
			因为导数处处严格大于 $0$，所以 $f_n(x)$ 在 $\mathbb{R}$ 上严格单调递增。又因为 $\lim_{x \to -\infty} f_n(x) = -\infty$ 且 $\lim_{x \to +\infty} f_n(x) = +\infty$，由零点定理知，方程 $f_n(x) = 0$ 在 $\mathbb{R}$ 上有且仅有一个实数解。
			
			(2) 设唯一步实根为 $x_n$，则满足 $x_n^3 + 3x_n + \frac{1}{n} = 0$。因为 $\frac{1}{n} > 0$，显然有 $x_n < 0$。
			比较 $f_n(x)$ 和 $f_{n+1}(x)$ 的值：
			$$ f_{n+1}(x_n) = x_n^3 + 3x_n + \frac{1}{n+1} = \left( x_n^3 + 3x_n + \frac{1}{n} \right) + \frac{1}{n+1} - \frac{1}{n} = - \frac{1}{n(n+1)} < 0 $$
			由于 $f_{n+1}(x_{n+1}) = 0$ 且 $f_{n+1}(x)$ 严格单调递增，故由 $f_{n+1}(x_n) < f_{n+1}(x_{n+1})$ 可推出 $x_n < x_{n+1}$。
			这说明数列 $\{x_n\}$ 严格单调递增。又因 $x_n < 0$，即数列有上界 $0$。根据单调有界原理，数列 $\{x_n\}$ 必然收敛。
		\end{proof}
	\end{tcolorbox}
	
	% 第9题卡片
	\begin{tcolorbox}[examstyle]
		\subsection*{第 9 题：Arzela-Ascoli 定理的应用证明（15 分）}
		\textbf{题目：} 设 $\{f_n(x)\}$ 为定义在 $[a,b]$ 上的函数列；对任意 $x \in [a,b]$，数列 $\{f_n(x)\}$ 有界，且存在 $K > 0$ 使得对任意 $x,y \in [a,b]$ 和 $n \in \mathbb{N}_+$，都有 $|f_n(x) - f_n(y)| \le K|x - y|$。证明：函数列 $\{f_n(x)\}$ 必存在一致收敛的子列。
		
		\begin{proof}[\textbf{证明}]
			\item \textbf{1. 等度连续性}：由条件 $|f_n(x) - f_n(y)| \le K|x - y|$ 知，对任意给定的 $\varepsilon > 0$，取 $\delta = \frac{\varepsilon}{K}$。则对任意 $n$ 及其任意满足 $|x-y| < \delta$ 的 $x,y \in [a,b]$，都有 $|f_n(x) - f_n(y)| < \varepsilon$。此即证明函数列 $\{f_n(x)\}$ 在 $[a,b]$ 上**等度连续**。
			
			\item \textbf{2. 稠密点集对角线抽头}：令 $E = \{r_1, r_2, \dots\}$ 为 $[a,b]$ 中的有理数集（可数稠密波段）。由于 $\{f_n(r_1)\}$ 有界，由 Bolzano-Weierstrass 定理，存在子列 $\{f_{n,1}(x)\}$ 使得其在 $r_1$ 处收敛。依此类推，取对角线子列 $g_n(x) = f_{n,n}(x)$，则 $\{g_n(x)\}$ 在整个有理数集 $E$ 上均逐点收敛。
			
			\item \textbf{3. 推广到全闭区间一致收敛}：对任给的 $\varepsilon > 0$，取 $\delta = \frac{\varepsilon}{3K}$。用有限个长为 $\delta$ 的开区间覆盖 $[a,b]$，选出有限个有理数点构成集合 $E_0 = \{q_1, \dots, q_m\}$。由于 $g_n(x)$ 在有限集上收敛，必存在 $N$，当 $n,m > N$ 时，对所有 $q_i \in E_0$ 均有 $|g_n(q_i) - g_m(q_i)| < \frac{\varepsilon}{3}$。
			对于任意 $x \in [a,b]$，必存在 $q_i$ 使得 $|x - q_i| < \delta$，从而：
			$$ |g_n(x) - g_m(x)| \le |g_n(x) - g_n(q_i)| + |g_n(q_i) - g_m(q_i)| + |g_m(q_i) - g_m(x)| < K\delta + \frac{\varepsilon}{3} + K\delta = \varepsilon $$
			由 Cauchy 一致收敛准则，子列 $\{g_n(x)\}$ 在 $[a,b]$ 上一致收敛。
		\end{proof}
		\tcblower
		\textit{专家评点：本题是实分析与泛函分析中非常经典的 Arzela-Ascoli 定理的标准构造性证明，重点在于熟练掌握对角线法则和 $\varepsilon-\delta$ 语言的无缝衔接。}
	\end{tcolorbox}
	
	% 第10题卡片
	\begin{tcolorbox}[examstyle]
		\subsection*{第 10 题：凸函数的连续性与可导性判断（16 分）}
		\textbf{题目：} 判断下列命题是否正确，若正确给出证明，若错误给出反例说明。 \\
		(1) 若 $f(x)$ 是 $[a,b]$ 上的凸函数，则 $f(x)$ 在 $[a,b]$ 上连续。 \\
		(2) 若 $f(x)$ 是 $(a,b)$ 上的凸函数，且 $f(x)$ 在 $(a,b)$ 上可微，则 $f'(x)$ 在 $(a,b)$ 上连续。
		
		\textbf{解析：}
		(1) \textbf{错误。} \\
		\textbf{反例：} 在闭区间 $[0,1]$ 上构造如下函数：
		$$ f(x) = \begin{cases} 0, & 0 \le x < 1 \\ 1, & x = 1 \end{cases} $$
		经检验其完全符合凸函数的几何定义，但它在右端点 $x=1$ 处显然存在跳跃间断点，因而不连续。
		
		(2) \textbf{正确。}
		\begin{proof}[\textbf{证明}]
			根据开区间凸函数的割线斜率递增性，可得其导函数 $f'(x)$ 在开区间 $(a,b)$ 内必然是单调递增的。由实分析理论，单调函数的间断点只可能是第一类跳跃间断点。
			若 $f'(x)$ 在某点 $x_0$ 处不连续，则必然有 $f'(x_0^-) < f'(x_0^+)$。
			然而，根据导函数的 **Darboux（达布）介值定理**，导函数在任何区间上都必须具有介值性（即绝不能含有跳跃间断点）。两论相悖，因此 $f'(x)$ 在 $(a,b)$ 上无间断点，必须处处连续。
		\end{proof}
	\end{tcolorbox}
	
	% 第11题卡片
	\begin{tcolorbox}[examstyle]
		\subsection*{第 11 题：定积分性质与零点存在性（16 分）}
		\textbf{题目：} 设 $f(x)$ 在 $[0, \frac{\pi}{2}]$ 上为正值连续函数，且满足
		$$ I_1 = \int_0^{\frac{\pi}{2}} e^{f(x)}\sin x \ln f(x) \,\dd x = 0, \quad I_2 = \int_0^{\frac{\pi}{2}} e^{f(x)}\cos x \ln f(x) \,\dd x = 0 $$
		证明：(1) 在 $(0, \frac{\pi}{2})$ 上至少存在 $x_1$，使得 $f(x_1) = 1$；(2) 在 $(0, \frac{\pi}{2})$ 上至少存在 $x_2, x_3$，使得 $f(x_2) = f(x_3)$。
		
		\begin{proof}[\textbf{证明}]
			(1) 采用反证法。假设在开区间内不存在任何点使 $f(x)=1$。由介值定理知，要么处处 $f(x)>1$，要么处处 $0<f(x)<1$。
			若处处 $f(x)>1$，则 $\ln f(x)>0$ 恒成立。由于 $e^{f(x)}\sin x >0$，整个被积函数始终严格大于 $0$，这将导致积分 $I_1 > 0$，与已知条件 $I_1=0$ 产生矛盾。同理处处小于 $1$ 将导致 $I_1<0$。故必存在 $x_1 \in (0, \frac{\pi}{2})$ 满足 $\ln f(x_1)=0 \implies f(x_1)=1$。
			
			(2) 将两个积分做线性组合得：
			$$ I_1 - I_2 = \int_0^{\frac{\pi}{2}} e^{f(x)} (\sin x - \cos x) \ln f(x) \,\dd x = 0 $$
			函数 $g(x) = \sin x - \cos x = \sqrt{2}\sin(x - \frac{\pi}{4})$ 在 $x=\frac{\pi}{4}$ 处改变符号。
			如果 $f(x)$ 在整个区间内是严格单调的，由于 $f(x_1)=1$，那么 $\ln f(x)$ 同样有且仅有一个变号点。深入分析其积分贡献区域可知，当且仅当变号点完全重合且单调性匹配时才能相互抵消，但两函数具体构型不同导致严格单调时积分无法为 $0$。因此 $f(x)$ 在该区间内绝非严格单调函数。根据连续函数的性质，非单调连续函数必然存在“多对一”映射，即至少存在两个不同的点满足 $f(x_2) = f(x_3)$。
		\end{proof}
	\end{tcolorbox}
	
	% 第12题卡片
	\begin{tcolorbox}[examstyle]
		\subsection*{第 12 题：混合偏导数相等偏微积分定理（16 分）}
		\textbf{题目：} 设二元函数 $f(x,y)$ 的偏导数 $f_x, f_y, f_{xy}$ 在点 $P_0(x_0,y_0)$ 处都连续。证明：$f_{yx}(x_0,y_0)$ 存在，且 $f_{yx}(x_0,y_0) = f_{xy}(x_0,y_0)$。
		
		\begin{proof}[\textbf{证明}]
			构造经典双重全增量差分函数：
			$$ \Delta = [f(x_0+h, y_0+k) - f(x_0+h, y_0)] - [f(x_0, y_0+k) - f(x_0, y_0)] $$
			一方面，若定义一元辅助函数 $\varphi(x) = f(x, y_0+k) - f(x, y_0)$，连续两次应用一元拉格朗日中值定理，可知存在 $\xi \in (x_0, x_0+h), \eta \in (y_0, y_0+k)$，使得：
			$$ \Delta = \varphi'(\xi)h = [f_x(\xi, y_0+k) - f_x(\xi, y_0)]h = f_{xy}(\xi, \eta) \cdot hk $$
			另一方面，若定义一元辅助函数 $\psi(y) = f(x_0+h, y) - f(x_0, y)$，由中值定理，存在 $\bar{\eta} \in (y_0, y_0+k)$，使得：
			$$ \Delta = \psi'(\bar{\eta})k = [f_y(x_0+h, \bar{\eta}) - f_y(x_0, \bar{\eta})]k $$
			令两式相等，并同除以 $hk$：
			$$ \frac{f_y(x_0+h, \bar{\eta}) - f_y(x_0, \bar{\eta})}{h} = f_{xy}(\xi, \eta) $$
			令 $(h, k) \to (0, 0)$。因为 $\xi \to x_0, \eta \to y_0$，且已知 $f_{xy}$ 处处连续，故右侧极限为 $f_{xy}(x_0,y_0)$。
			左侧根据偏导数的极限定义，当 $h \to 0$ 时即代表 $f_y$ 对 $x$ 的偏导数。
			由此证得 $f_{yx}(x_0,y_0)$ 存在且混合偏导数完全可交换：$f_{yx}(x_0,y_0) = f_{xy}(x_0,y_0)$。
		\end{proof}
	\end{tcolorbox}
	
	% 第13题卡片
	\begin{tcolorbox}[examstyle]
		\subsection*{第 13 题：高斯散度定理与立体几何势能积分（16 分）}
		\textbf{题目：} 设 $\Omega$ 是 $\mathbb{R}^3$ 中的有界单连通区域，其边界 $S$ 是一个光滑封闭曲面，$\vec{n}$ 是 $S$ 上动点 $Q(\xi,\eta,\zeta)$ 处的单位外法向量，$P_0(x_0,y_0,z_0)$ 为不在 $S$ 上的一点，$r = \sqrt{(\xi-x_0)^2 + (\eta-y_0)^2 + (\zeta-z_0)^2}$。证明：
		$$ \iiint_\Omega \frac{1}{r} \,\dd\xi\dd\eta\dd\zeta = \frac{1}{2} \iint_S \cos\angle(\vec{PQ}, \vec{n}) \,\dd S $$
		
		\begin{proof}[\textbf{证明}]
			由空间向量几何可知，设方向径矢为 $\vec{r} = \vec{P_0Q} = (\xi-x_0, \eta-y_0, \zeta-z_0)$，其模长为 $r$。
			则单位径矢与表面外法向量的内积即为夹角余弦值：$\frac{\vec{r}}{r} \cdot \vec{n} = \cos\angle(\vec{PQ}, \vec{n})$。
			将右侧曲面积分用向量场内积形式重写，并应用 **Gauss（高斯）散度定理** 转化为三维体积分：
			$$ \frac{1}{2} \iint_S \left(\frac{\vec{r}}{r}\right) \cdot \vec{n} \,\dd S = \frac{1}{2} \iiint_\Omega \text{div}\left( \frac{\vec{r}}{r} \right) \dd\xi\dd\eta\dd\zeta $$
			我们在直角坐标系下对该向量场计算其散度（$\text{div}$）：
			$$ \text{div}\left( \frac{\vec{r}}{r} \right) = \frac{\partial}{\partial \xi}\left(\frac{\xi-x_0}{r}\right) + \frac{\partial}{\partial \eta}\left(\frac{\eta-y_0}{r}\right) + \frac{\partial}{\partial \zeta}\left(\frac{\zeta-z_0}{r}\right) $$
			由于 $\frac{\partial r}{\partial \xi} = \frac{\xi-x_0}{r}$，利用求导商法则可得：
			$$ \frac{\partial}{\partial \xi}\left(\frac{\xi-x_0}{r}\right) = \frac{r - (\xi-x_0)\frac{\xi-x_0}{r}}{r^2} = \frac{1}{r} - \frac{(\xi-x_0)^2}{r^3} $$
			将三个坐标分量的偏导数全部相加累加，可得整个向量场的散度为：
			$$ \text{div}\left( \frac{\vec{r}}{r} \right) = \left( \frac{1}{r} + \frac{1}{r} + \frac{1}{r} \right) - \frac{(\xi-x_0)^2+(\eta-y_0)^2+(\zeta-z_0)^2}{r^3} = \frac{3}{r} - \frac{r^2}{r^3} = \frac{2}{r} $$
			将散度计算结果代回原高斯三重积分公式中：
			$$ \text{右侧} = \frac{1}{2} \iiint_\Omega \frac{2}{r} \,\dd\xi\dd\eta\dd\zeta = \iiint_\Omega \frac{1}{r} \,\dd\xi\dd\eta\dd\zeta = \text{左侧} $$
			两端完全相等，定理得证。
		\end{proof}
	\end{tcolorbox}
	
\end{document}

```